% !TeX root = ../main.tex
% Language revision R20261107-10; numerical and inferential scope retained.
\section{Discussion and Conclusion}
\subsection{Findings and Contributions}

For RQ1, adjustment for era composition reduced the friendly--other-match gap, and geographic comparisons depended on the opponents represented. Several model-conditional ranges were too broad to resolve the direction of a difference. These comparisons provide a common descriptive basis for the strength-adjusted analysis, without identifying a pure home effect or a universal regional ranking.

For RQ2, the larger reduction in aggregate loss was associated with strength, while context provided a smaller increment. Dynamic paired comparisons~\citep{glickman1999dynamic} convert earlier results into estimates of ability and uncertainty; the common-sample comparison then tests whether context improves predictions beyond those estimates. The result gives context a conditional forecasting role. It does not turn geographic or tournament differences in observed home-win rates into an independent measure of venue advantage.

For RQ3, cumulative history performed best in the main comparison among the tested choices and remains the default. A shorter outcome-fitting window reduced computation at a small predictive cost. Regional diagnostics qualified the aggregate context gain: in Africa, its direction was sensitive to particular years and teams, while opposing probability biases persisted in selected uncertainty groups. These findings identify specific targets for further evaluation and calibration.

The project's contribution is the integration of match and identity curation, a date-ordered strength-estimation and prediction implementation, and common-sample evaluation of contextual information. The paired-comparison, approximate-inference, calibration, and scoring principles are drawn from the cited literature; methodological integration and empirical analysis, rather than the invention of those principles, distinguish this work. Documented information dates and an explicit withholding rule connect the historical analysis to prospective application. Proper probability scores~\citep{gneiting2007scoring} identify differences that mean bias or accuracy alone may obscure. The archived release permits subsequent evaluation against outcomes unavailable at issue time.

\subsection{Limitations}

The findings apply to the documented study population. Some historical identity and membership decisions rely on records accepted by the project lead, and the full archive has not been independently verified. Geographic classifications include study-specific overrides, while sparse opponent combinations limit comparability. Because some scores include extra time, the target cannot consistently be interpreted as a regulation-time result.

Approximate strength inference, assumed jump dynamics, and restricted context terms may leave substantial variation unexplained. The conditional comparisons do not account for all differences in travel, team selection, venues, or match importance. The main model combines draws and away wins, whereas the auxiliary model distinguishes three outcomes; their classification accuracies address different targets.

Some limitations have measurable consequences. The ordinary and strict numerical replays differed little, whereas initialization alternatives changed raw probabilities by more than $1$~pp in 2,905 study matches. These changes measure sensitivity and do not provide additional confidence bounds for corrected predictions. The final pipeline's recent-period loss was also slightly higher than that of the fixed-step baseline, limiting claims about the benefit of persistent jumps and probability correction. The separate effects of Gaussian projection, unmeasured squad changes, weak cross-region connections, and historical identity errors have not been quantified.

Development results changed both the integration procedure and probability correction: numerical-reference failures prompted a new integration order, and residual bias prompted separate shape and drift corrections. These revisions addressed observed defects under the documented criteria. Because the same archive also informed those decisions, the resulting historical scores can favor revisions specific to that archive. The prospective release preserves predictions made before their outcomes and permits a later test on new matches.

Chronological prediction excludes each match's own outcome from its inputs, but it does not make a reused development archive an untouched final test. Simulation results apply to their specified generators, and no dependence-adjusted interval from real match data establishes the persistence of the small contextual gain. Deletion checks assess evaluation composition while retaining the original learning history. Prospective evaluation and regional diagnostics therefore address different remaining questions about model use.

The historical uncertainty procedure also failed under a permanent global-shock generator. Its envelopes therefore remain conditional on the generating mechanisms for which the procedure met the validation criteria, rather than providing guarantees for arbitrary football histories.

\subsection{Future Work}

Four extensions warrant investigation. First, a three-class joint model could use a draw interval, as in TrueSkill~\citep{herbrich2006trueskill}, or divide non-home-win probability between draws and away wins. Both formulations require evaluation against the existing auxiliary model.

Second, calibration mappings learned from earlier predictions and outcomes could reduce systematic bias. Temperature scaling~\citep{guo2017calibration} is one candidate, although its published evidence concerns neural classifiers. Benefits require chronological and regional evaluation, and improved calibration need not improve classification accuracy.

Third, goal-based models~\citep{dixon1997scores} could estimate richer outcomes, subject to consistent score definitions and joint treatment of scores and match results. Verified venue or squad information must be available before prediction. These variables remain untested explanations of the observed biases.

Fourth, club-football analysis requires new ratings and parameters for seasons, transfers, promotion and relegation, and connections between leagues. An initial single-league study could model seasonal transitions explicitly. The predictive benefits of these proposed extensions remain unestablished.

\subsection{Conclusion}

Within this international-football archive, strength provided the larger reduction in predictive loss, and context supplied a smaller increment whose stability varied across regional checks. Historical-window comparisons favored cumulative history among the tested choices and identified a computation--prediction trade-off. These results support evaluating historical home-win patterns as additional forecasting information after accounting for uncertain strength. The archived prospective release applies that approach with explicit information dates and preserves its probabilities for evaluation when outcomes become available.

\section*{Declarations}
\phantomsection\addcontentsline{toc}{section}{Declarations}
\paragraph{Data and code availability.}
Source match data are available under CC0 from martj42~\citep{martj42_data}. The accompanying repository provides code, configurations, build instructions, experiment records, and input/output fingerprints identifying the reported versions.

The English guide \path{paper/SUPPORTING_RESULTS.md} indexes all stage tables, methods, reproduction instructions, windows, geographic matchups, sensitivity outputs, and the 52-fixture forecast ledger. The main text reports selected results.

\paragraph{Team responsibilities.}
The division of responsibilities was as follows: Ruiyi Gao was responsible for study design, coordination, computational analysis, and manuscript preparation; Zhenglin Chen for code review and methodological refinement; and Wenyu Jiang for supplementary data collection and presentation preparation.

\paragraph{Use of AI-assisted tools.}
AI-assisted tools supported code development, documentation, and manuscript revision. Empirical claims are supported by the archived data and computational outputs, rather than by generated text alone.
